\documentclass[11pt]{article}
\usepackage[margin=1in]{geometry}
\usepackage{amsmath,amssymb,microtype}
\usepackage[colorlinks=true,urlcolor=blue,citecolor=blue]{hyperref}
\newcommand{\Mn}{M_n(\mathbb C)}
\title{A negative literature answer to the tracial Markov-map question in arXiv:0803.2979}
\author{Literature-status note for \texttt{fa\_banach\_001}}
\date{11 August 2026}

\begin{document}
\maketitle

\noindent\textbf{Status.}
\texttt{literature\_already\_answered (negative)}.

\section*{Original question}

After constructing factorizations for real positive Schur multipliers,
Ricard observes that compositions in different bases, conditional
expectations, representations, and convex combinations yield a wide class
of factorizable maps.  He then asks on page 4 of the arXiv PDF
\cite{Ricard2008} whether these constructions can ``achieve all tracial
markovian maps for $M_n$.''

Every map produced in this way is factorizable: this is the point of the
preceding construction and its permanence properties.  Consequently, any
tracial Markov map that is not factorizable is a counterexample to the
question.

\section*{Explicit negative answer}

Haagerup and Musat give exactly such a map in Example 3.1 on pages 9--10 of
arXiv:1009.0778 \cite{HaagerupMusat2011}.  Put
\[
 a_1=\frac1{\sqrt2}\begin{pmatrix}0&0&0\\0&0&-1\\0&1&0\end{pmatrix},\quad
 a_2=\frac1{\sqrt2}\begin{pmatrix}0&0&1\\0&0&0\\-1&0&0\end{pmatrix},\quad
 a_3=\frac1{\sqrt2}\begin{pmatrix}0&-1&0\\1&0&0\\0&0&0\end{pmatrix}
\]
and define
\[
 T(x)=\sum_{i=1}^3a_i^*xa_i,\qquad x\in M_3(\mathbb C).
\]
Since
\[
 \sum_i a_i^*a_i=\sum_i a_i a_i^*=I_3,
\]
$T$ is completely positive, unital, and trace preserving, hence is a
$\tau_3$-Markov map.

The same example proves that $T$ is not factorizable.  Indeed, their
factorization criterion would give elements $v_1,v_2,v_3$ in a finite von
Neumann algebra such that $u=\sum_i a_i\otimes v_i$ is unitary.  The diagonal
relations in $u^*u=I$ force $v_i^*v_i=I$, hence finiteness makes every $v_i$
unitary.  The off-diagonal relations simultaneously force
\[
 v_1^*v_2=v_2^*v_3=v_3^*v_1=0,
\]
an impossibility for unitaries.  Thus $T$ is a tracial Markov map on $M_3$
outside the factorizable class, and so it cannot arise from Ricard's listed
construction.  This answers the question negatively.

\section*{Dimension scope and identification}

The introduction of the answering paper explicitly cites Ricard's paper as
the source for factorizability of real positive Schur multipliers and then
announces nonfactorizable $\tau_n$-Markov maps for every $n\ge3$.  It also
records K\"ummerer's positive result that every trace-preserving Markov map
on $M_2(\mathbb C)$ is factorizable.  Hence the later literature supplies the
sharp dimension boundary: the exhaustion statement holds at $n=2$ but
fails for every $n\ge3$.

The original question was checked on source PDF page 4.  The counterexample
and its contradiction proof were checked on supporting PDF pages 9--10; the
all-$n\ge3$ statement is in the abstract and Remark 4.8.  This packet makes
no new mathematical claim.

\begin{thebibliography}{9}
\bibitem{Ricard2008}
É.~Ricard,
\emph{A Markov dilation for self-adjoint Schur multipliers},
Proc. Amer. Math. Soc. \textbf{136} (2008), 4365--4372;
arXiv:0803.2979,
\url{https://arxiv.org/abs/0803.2979}.

\bibitem{HaagerupMusat2011}
U.~Haagerup and M.~Musat,
\emph{Factorization and dilation problems for completely positive maps on
von Neumann algebras},
Comm. Math. Phys. \textbf{303} (2011), 555--594;
arXiv:1009.0778,
\url{https://arxiv.org/abs/1009.0778}.
\end{thebibliography}

\end{document}

